\documentclass[11pt]{article}

\usepackage[T1]{fontenc}
\usepackage{lmodern}
\usepackage{microtype}
\usepackage{amsmath,amssymb,amsthm}
\usepackage{graphicx}
\usepackage[margin=1in]{geometry}
\usepackage[dvipsnames]{xcolor}
\usepackage[colorlinks=true,linkcolor=MidnightBlue,citecolor=MidnightBlue,urlcolor=MidnightBlue]{hyperref}

\definecolor{navy}{HTML}{17324D}
\definecolor{pale}{HTML}{EEF4F8}
\newtheorem{theorem}{Theorem}
\newcommand{\Fin}{\mathrm{Fin}}
\newcommand{\LGBPs}{\mathrm{(LGBP_s)}}
\newcommand{\LGBPl}{\mathrm{(LGBP_l)}}
\newcommand{\cJ}{\mathcal{J}}
\newcommand{\cP}{\mathcal{P}}

\title{A full answer to Question 8.4:\\
\(\Fin\otimes\Fin\) separates \(\LGBPl\) from \(\LGBPs\)}
\author{Automated Mathematics Run \texttt{fa\_banach\_001}}
\date{19 July 2026}

\begin{document}
\maketitle

\noindent\colorbox{pale}{\parbox{0.96\textwidth}{
\textbf{Status.} Candidate full solution, likely valid; independent human
review pending.}}

\section{Source question}

Sobota and \.{Z}uchowski ask the following in
Question~8.4 on PDF page~47 of their preprint~\cite{SZ}.

\begin{center}
\includegraphics[width=0.98\textwidth]{figures/open_problem_crop.png}
\end{center}

In words: does there exist an ideal on \(\omega\) which has the
Local-to-Global Boundedness Property for lower semicontinuous submeasures
\(\LGBPl\), but not the corresponding property for arbitrary submeasures
\(\LGBPs\)?

\section{Definitions}

Let \(\Omega=\omega\times\omega\), and write
\(R_n=\{n\}\times\omega\) and
\(A_n=\{k:(n,k)\in A\}\).  We work with the standard ideal
\[
 \cJ=\Fin\otimes\Fin
   =\bigl\{A\subseteq\Omega:\{n:A_n\text{ is infinite}\}\text{ is finite}\bigr\}.
\]
Since \(\Omega\) is countable, any bijection \(\omega\to\Omega\) transports
\(\cJ\) to an ideal on \(\omega\).

A submeasure is a monotone, subadditive map
\(\varphi:\cP(\Omega)\to[0,\infty]\) which vanishes at the empty set and is
finite on singletons.  Put
\(\Fin(\varphi)=\{A:\varphi(A)<\infty\}\).  An ideal \(\mathcal I\) has
\(\LGBPs\) if every submeasure \(\varphi\) with
\(\mathcal I\subseteq\Fin(\varphi)\) is uniformly bounded on \(\mathcal I\).
It has \(\LGBPl\) if every \emph{lower semicontinuous} submeasure \(\psi\)
with \(\mathcal I\subseteq\Fin(\psi)\) satisfies
\(\psi(\Omega)<\infty\)~\cite{SZ}.

\section{Proof intuition}

Counting the infinite rows gives exactly the non-lower-semicontinuous
submeasure needed to destroy \(\LGBPs\).  Lower semicontinuity prevents this
phenomenon: if an lsc submeasure were infinite on the whole grid but finite on
\(\cJ\), then every tail of rows would still have infinite value.  One can
choose a finite set of arbitrarily large value from each successive tail.
Their diagonal union remains finite in every row, hence belongs to \(\cJ\),
but has infinite value---the required contradiction.

\begin{theorem}
The ideal \(\cJ=\Fin\otimes\Fin\) has \(\LGBPl\) and does not have
\(\LGBPs\).  Consequently, Question~8.4 has an affirmative answer in ZFC.
\end{theorem}

\begin{proof}
For \(A\subseteq\Omega\), set
\[
 S(A)=\{n\in\omega:A_n\text{ is infinite}\},\qquad
 \rho(A)=|S(A)|\in\mathbb N\cup\{\infty\}.
\]
This is a submeasure.  Indeed, it is monotone, vanishes on finite sets, and
\[
 S(A\cup B)\subseteq S(A)\cup S(B),
\]
because the union of two finite row sections is finite.  Thus \(\rho\) is
subadditive.  Moreover,
\(\Fin(\rho)=\cJ\), while for
\(P_m=\bigcup_{n<m}R_n\) one has \(P_m\in\cJ\) and \(\rho(P_m)=m\).
Hence \(\rho\) is finite but unbounded on \(\cJ\), proving that \(\cJ\)
does not have \(\LGBPs\).

Now let \(\psi\) be an lsc submeasure such that
\(\cJ\subseteq\Fin(\psi)\).  Suppose for a contradiction that
\(\psi(\Omega)=\infty\).  For every \(m\), the initial row block \(P_m\)
belongs to \(\cJ\), so \(\psi(P_m)<\infty\).  With
\[
 T_m=\Omega\setminus P_m=\bigcup_{n\ge m}R_n,
\]
subadditivity forces \(\psi(T_m)=\infty\); otherwise
\(\psi(\Omega)\le\psi(P_m)+\psi(T_m)<\infty\).

By lower semicontinuity, for every \(m\) there is a finite set
\(F_m\subseteq T_m\) such that \(\psi(F_m)>m\).  Let
\(A=\bigcup_m F_m\).  Fix a row \(R_n\).  Since
\(F_m\subseteq T_m\), the set \(F_m\) misses \(R_n\) whenever \(m>n\).
Therefore
\[
 A\cap R_n\subseteq\bigcup_{m\le n}F_m
\]
is finite.  Every row section of \(A\) is finite, so \(A\in\cJ\).  On the
other hand, monotonicity gives
\(\psi(A)\ge\psi(F_m)>m\) for every \(m\), hence \(\psi(A)=\infty\),
contradicting \(\cJ\subseteq\Fin(\psi)\).  Thus
\(\psi(\Omega)<\infty\), which is precisely \(\LGBPl\).
\end{proof}

\section{Checks, scope, and novelty}

The witness \(\rho\) is not lower semicontinuous: finite approximations to a
full row have value zero, while the full row has value one.  Hence it does not
conflict with the \(\LGBPl\) conclusion.  Also, \(\cJ\) is not a P-ideal:
the full rows \(R_n\) have no pseudounion in \(\cJ\).  The example therefore
lies outside the analytic P-ideal class singled out immediately before the
source question.

The run's lightweight indexes and a bounded web/arXiv search were queried by
the source id, the exact LGBP terminology, \(\Fin\otimes\Fin\), and close
variants.  The search found the source and general literature on this ideal,
but no later answer to Question~8.4 or matching separation argument.  This is
not an exhaustive novelty claim.

\paragraph{Human-review focus.}
Check the two LGBP quantifiers, the lsc selection of the finite sets \(F_m\),
and the row-finiteness of their diagonal union.

\begin{thebibliography}{9}
\bibitem{SZ}
D. Sobota and T. \.{Z}uchowski,
\emph{The Nikodym and Grothendieck properties of Boolean algebras and rings
related to ideals}, arXiv:2510.19744 (2025), Question~8.4.
\end{thebibliography}

\end{document}
